% Section 7.1, from lectures/L07/README.md: the objectives, the evaluation questions, and the next
% lecture.

\needspace{10\baselineskip}
\section{Review}
\label{c7:sec:review}

\subsection*{What you should be able to do now}
After this chapter, you should:
\begin{itemize}
  \item Be able to implement a simple conv, max pooling, and flatten layer in software.
\end{itemize}

\needspace{6\baselineskip}
\subsection*{Questions to test yourself}
You should be able to explain:
\begin{enumerate}
  \item Which parts of your \code{ml::ConvLayer} implementation correspond to feedforward and
    backpropagation, respectively?
  \item Why do we use the same kernel (the same weights) across the entire image instead of unique
    weights per position?
  \item How does your \code{ml::MaxPoolLayer} implementation know which positions the gradients
    should be propagated back to during backpropagation?
  \item Why is a flatten layer needed between the convolutional and pooling layers and the dense
    layer?
\end{enumerate}

\needspace{6\baselineskip}
\subsection*{Looking ahead}
Chapter~\ref{c8:ch:conv} implements the conv layer in the real, polymorphic CNN pipeline,
\file{cnn_work}, replacing the stub it starts out with.
